% ---------------------------------------------------------------------------
% Shared section: A3 data handling (written by Kolipakula Sai Tushaar).
% DONE - no edits needed. Compiled through report/main.tex.
% ---------------------------------------------------------------------------
\section{A3 Data Handling}
\label{sec:a3-prep}

A3 starts from \texttt{master\_cleaned.csv}, the output of the A1 pipeline
(80{,}000 rows, 43 columns, no nulls or duplicates; see the A1 report in the Appendix). The script
\texttt{preprocess\_a3.py} adds eight steps that the A3 models need and that A1
did not. Every step asserts its own precondition, so the script fails loudly
if the input ever changes. It writes one shared table, \texttt{a3\_base.csv.gz},
which all three members load through \texttt{a3\_common.load\_base()} so that
every analysis uses the same features, split and colours.

\begin{enumerate}
  \item \textbf{Re-validation.} Schema, null, duplicate and range checks on 14
        numeric columns are repeated before any modelling.
  \item \textbf{Exact redundancies removed.} Three columns are formulas of
        others for every row and would double-weight one variable in a
        distance-based or linear model:
        \begin{itemize}
          \item \texttt{exam\_anxiety\_score} $=\mathrm{clip}(15-\text{motivation},5,10)$;
          \item \texttt{study\_hours\_per\_week} $=7\times$ daily study hours;
          \item \texttt{leisure\_screen\_hours} $=$ social media $+$ streaming hours.
        \end{itemize}
  \item \textbf{Leakage recorded.} \texttt{dropout\_risk} equals
        $(\text{stress}>8)\wedge(\text{motivation}\le 3)$ for all rows
        (A1 report, Fig.~21), and \texttt{performance\_band} and
        \texttt{performance\_delta} are functions of the exam score. These
        rules are stored in \texttt{a3\_feature\_groups.json} so no model uses
        them as features by mistake.
  \item \textbf{Ceiling flags.} 25.8\% of students score exactly 100 and
        40.0\% have a prior GPA of exactly 4.0; 21.4\% are at both caps. Both
        variables are censored, which matters for any residual or linear model.
  \item \textbf{Prior-attainment bands.} \texttt{gpa\_band} on
        \texttt{previous\_gpa\_pct}: Below 80 (17{,}487 students, mean score
        71.7), 80--90 (15{,}154; 85.2), 90--99 (15{,}390; 94.4) and Elite, GPA
        exactly 4.0 (31{,}969; 98.0).
  \item \textbf{Outlier screen.} The $1.5\times$IQR rule flags 0.25--0.45\% of
        values in five columns. All lie inside the instrument ranges, so they
        are kept as genuine tail values.
  \item \textbf{Encoding.} Booleans become 0/1; four nominal columns (gender,
        major, study environment, learning style) become 18 one-hot columns;
        ordinal columns keep their A1 integer codes.
  \item \textbf{Train/test split.} 80/20, stratified on performance band
        $\times$ dropout flag, seed 42. Train: 64{,}000 rows, mean score 89.13,
        dropout rate 1.98\%; test: 16{,}000 rows, 89.19, 1.98\%.
\end{enumerate}

\paragraph{Supplementary data.} None is used. Every question in A3 can be
answered from the A1 dataset itself, and the strongest explanatory variable,
prior GPA, is already in it. The VAST 2026 Challenge material mentioned in the
brief was not used.

Scaling is deliberately left out of the shared file. Each member fits the
scaler on the training split inside their own pipeline, so no information from
the test split reaches any model.
